% I. Анализ на предметната област и литературен преглед.
% Тук се прави връзката между учебното съдържание на дисциплината и обхвата на проекта.
\section{Анализ на предметната област и литературен преглед}
\label{sec:literature}

Дисциплината „Мобилни и стационарни комуникационни мрежи“ обхваща два слоя, които се
изучават заедно, но се проектират по различни принципи. Първият е преносната среда със
своите разновидности: стационарни кабелни мрежи, безжични мрежи с локален обхват,
клетъчни мрежи с преминаване между клетки, спътникови връзки с голямо закъснение на
разпространение и мрежи с виртуални канали от типа ATM. Вторият е мрежовият слой на
стека TCP/IP, който работи над всяка от тези среди и трябва да остава функционален,
когато характеристиките им се променят. Учебното съдържание отделя специално внимание
на принципите на маршрутизация и управление на потоците в TCP/IP базирани мрежи и
именно този слой е предмет на настоящия проект.

Класификацията на протоколите за маршрутизация вътре в една автономна система е
установена и се свежда до две семейства. Протоколите по вектор на разстоянието обменят
със съседите си собствената си маршрутна таблица, като RIP~\cite{rfc1058}. Те са прости
за реализация, но страдат от бавно разпространение на лошите новини и изискват
допълнителни механизми срещу образуването на цикли. Протоколите по състояние на
връзките разпространяват описание на локалната свързаност до всички възли, така че всеки
възел изгражда еднакъв граф на топологията и изчислява маршрутите самостоятелно.
OSPF версия 2~\cite{rfc2328} е каноничното описание на този подход, а версия 3 го
пренася върху IPv6~\cite{rfc5340}. Сравнителен анализ на двете семейства, включително
причините за преходните цикли, е даден при Perlman~\cite{perlman2000}.

Изчисляването на маршрутите при протоколите по състояние на връзките се свежда до
търсене на дърво от най-кратки пътища с корен текущия възел, тоест до алгоритъма на
Дийкстра~\cite{dijkstra1959}. Изборът на алгоритъм е безспорен, но времето за сходимост
не се определя от него: изчислението е бързо в сравнение с интервалите, през които
възлите изобщо научават за промяната. Тези интервали се задават от таймерите за
откриване на съседи и от закъснението на лавинното разпространение на обявленията.
Затова литературата за ускоряване на сходимостта се съсредоточава върху по-бързо
откриване на отказа, например чрез двупосочно откриване на препращането
(BFD)~\cite{rfc5880}, а не върху оптимизация на самото изчисление.

Управлението на потоците в TCP/IP се разполага в друг слой, но е свързано със
сходимостта по наблюдаем начин. Прозорецът за претоварване на TCP се свива при загуба
на пакети~\cite{jacobson1988,rfc5681}, а точно преходните цикли при пренастройване на
маршрутите произвеждат загуба, която за подателя е неразличима от загуба поради
претоварване. Реакцията на транспортния слой при промяна в маршрутизацията е следствие
от поведението на мрежовия слой и това прави времето за сходимост величина с пряк ефект
върху пропускателната способност, а не вътрешна характеристика на протокола. Общата
рамка на двата слоя е изложена при Kurose и Ross~\cite{kurose2021}, а нормативното
описание на TCP е в~\cite{rfc9293}.

Инструментите за изследване на такива преходи също са установени. Симулаторите с
дискретни събития OMNeT++~\cite{varga2008} и ns-3~\cite{riley2010} предоставят зряла
среда с модели на среди и протоколи. Те обаче симулират собствени модели на протокола,
докато настоящият проект се стреми към друго свойство: една и съща реализация да бъде
изпълнима както над симулирано, така и над реално време.
\TODO{да се добави преглед на публикации, измерващи сходимостта на OSPF в реални мрежи,
след като бъде подбран конкретен набор}
